\section{Introducción: por qué los robots necesitan un modelo base}

Esta sección plantea primero el problema de todo el episodio. El juicio central de Chen Jianyu (陈建宇) es que, en el pasado, la robótica solía funcionar con «100 escenarios, 100 tareas, 100 modelos especializados», un enfoque que difícilmente escala. El éxito de los modelos de lenguaje de gran tamaño y de los modelos multimodales le mostró a la robótica otro camino: ¿se puede entrenar para los robots un foundation model más general, que generalice entre tareas, entre cuerpos de robot y entre escenarios?

\begin{importantbox}{Tesis central de este episodio}
VLA, es decir, Vision-Language-Action Model, no consiste solo en conectar un modelo de lenguaje a un robot, sino en integrar visión, lenguaje y acción en un mismo sistema de extremo a extremo y escalable. Un verdadero modelo base para robots debe pasar de «usar los modelos fundacionales existentes para hacer robótica» a «preentrenar modelos fundacionales para los robots».
\end{importantbox}

\begin{figure}[H]
\centering
\includegraphics[width=0.96\textwidth]{ch03-00693.jpg}
\caption{How robotic learning works now? El aprendizaje robótico tradicional suele organizarse en torno al ciclo cerrado de tarea, datos, modelo y ejecución. Fotograma de la proyección, aprox. 00:11:33.}
\end{figure}

\begin{knowledgebox}{Lectura de la figura: la presión para pasar de modelos especializados a modelos generales}
La figura no presenta un artículo concreto, sino el contexto de toda la clase: el aprendizaje robótico ha dependido durante mucho tiempo de tareas específicas, datos específicos y despliegues específicos. Chen Jianyu subraya que este enfoque no es lo bastante general para hogares, fábricas y robots humanoides, por lo que se necesita la vía de los modelos base.
\end{knowledgebox}

\subsection{Las dos etapas de la IA para robots}

Esta sección distingue con claridad las dos vías. La primera etapa es leveraging foundation models in robotics, es decir, tomar modelos fundacionales existentes, como LLM, VLM o Code LM, para sustituir algún módulo del sistema robótico. La segunda etapa es pretraining foundation models for robotics, es decir, entrenar directamente para los robots modelos capaces de producir acciones. La primera consiste en ensamblar y apoyarse en lo existente; solo la segunda constituye un verdadero modelo base para robots.

\begin{figure}[H]
\centering
\includegraphics[width=0.90\textwidth]{ch05-01297.jpg}
\caption{Foundation models for robotics: relación entre módulos, desde LLM y VLA hasta Actuation. Fotograma de la proyección, aprox. 00:21:37.}
\end{figure}

\begin{knowledgebox}{Lectura de la figura: los tres bloques del robot y su sustitución por modelos fundacionales}
La figura reúne Perception, Actuation y LLM. Tradicionalmente, un robot se divide en planificación, percepción y ejecución; la ola de los modelos grandes sustituyó primero a Planning, después a Perception, y solo al final intentó que Action también se integrara en un modelo unificado.
\end{knowledgebox}

\begin{knowledgebox}{Términos clave: VLA, VLM, LLM}
\noindent\begin{tabularx}{\textwidth}{p{0.18\textwidth}p{0.34\textwidth}X}
\toprule
Término & Significado & Función en esta clase \\
\midrule
LLM & Large Language Model, modelo de lenguaje de gran tamaño & Primero sustituye o refuerza la planificación del robot. \\
VLM & Vision-Language Model, modelo de visión y lenguaje & Se usa para la percepción, la comprensión de la escena y la retroalimentación visual. \\
VLA & Vision-Language-Action Model, modelo de visión, lenguaje y acción & Conecta directamente visión, lenguaje y acción en un control de extremo a extremo. \\
Robot Foundation Model & Modelo base para robots & Su objetivo es generalizar entre tareas, entre cuerpos de robot y entre escenarios. \\
\bottomrule
\end{tabularx}
\end{knowledgebox}

\subsection{Resumen de la sección}

La aparición de VLA es la pista fundamental del paso de la robótica desde el ensamblaje de módulos hacia los modelos base generales. El hilo conductor de toda la clase va de apoyarse en LLM/VLM, a RT-2/OpenVLA/HiRT/pi0/GR00T N1, y de ahí a diffusion policy, RDT y los futuros modelos unificados de comprensión y predicción.
